\chapter{Valutazione dei Sistemi di Information Retrieval}
\label{chap:evaluation}

\FloatBarrier
\section{Valutare un Sistema di Information Retrieval: Metodologia e Obiettivi}
\label{sec:metodologia-obiettivi-valutazione}

Un sistema di Information Retrieval (IR) \`e un'infrastruttura ingegneristica complessa progettata per connettere persone e bisogni informativi a collezioni estese di dati non strutturati~\cite{manning2008,buttcher2010}. La valutazione scientifica e industriale di tali sistemi costituisce una disciplina cardine dell'IR, affrontata storicamente a partire dagli esperimenti pionieristici di Cyril Cleverdon negli anni '60~\cite{cleverdon1967}.

\FloatBarrier
\subsection{Efficienza, Servizi Accessori ed Efficacia Intrinseca}
\label{subsec:tre-dimensioni-valutazione}

Valutare la qualit\`a complessiva di un motore di ricerca richiede di disaccoppiare tre dimensioni di indagine complementari:

\begin{enumerate}
    \item \textbf{Metriche di Efficienza Operativa e di Sistema:}
    \begin{itemize}
        \item \textit{Throughput di indicizzazione:} Velocit\`a di assorbimento e indicizzazione di una collezione documentale, misurata tipicamente in documenti elaborati per ora ($\text{doc}/\text{h}$).
        \item \textit{Indicizzazione incrementale:} Capacit\`a del sistema di aggiornare l'indice con nuovi flussi di dati (ad esempio l'inserimento di $+10\,000$ nuovi prodotti al giorno in un portale e-commerce) senza interrompere la disponibilit\`a operativa né degradare i tempi di accesso.
        \item \textit{Throughput e latenza di ricerca:} Tempo di risposta medio e percentili critici (95° e 99° percentile, tipicamente $< 50\text{--}100\text{ ms}$) unitamente all'impegno di CPU, memoria RAM e I/O per l'elaborazione contemporanea di decine di migliaia di query al secondo su collezioni di scala planetaria (decine di milioni o miliardi di documenti).
    \end{itemize}
    \item \textbf{Qualit\`a ed Efficacia dei Servizi Accessori:}
    \begin{itemize}
        \item Valutazione della precisione dei suggerimenti ortografici (\textit{query spelling correction} o \textit{``forse cercavi: ...''}), dell'espansione e riscrittura automatica delle query, della pertinenza delle raccomandazioni di prodotti o pagine correlate, e della leggibilit\`a e densit\`a informativa dei riassunti contestuali (\textit{result snippets}).
        \item Sebbene questi servizi accessori migliorino l'esperienza d'uso, essi non certificano di per sé la bont\`a fondamentale dell'algoritmo di ranking.
    \end{itemize}
    \item \textbf{Efficacia Intrinseca della Graduatoria (Ranking Intrinseco):}
    \begin{itemize}
        \item Misura la capacit\`a fondamentale del motore di ricerca di soddisfare l'utente, preparando per ciascuna interrogazione una pagina dei risultati (\textit{Search Engine Results Page}, SERP) che posizioni i documenti pertinenti ai vertici della graduatoria.
        \item L'obiettivo primario \`e massimizzare la soddisfazione complessiva dell'utente (\textit{User Happiness}).
    \end{itemize}
\end{enumerate}

La Figura~\ref{fig:architettura-valutazione} schematizza la duplice pipeline offline/online di un moderno motore di ricerca, evidenziando i punti di rilevazione delle metriche di efficienza e di efficacia.

\begin{figure}[H]
\centering
\resizebox{0.98\linewidth}{!}{%
\begin{tikzpicture}[
    box/.style={rectangle, rounded corners=4pt, draw=navy!85!black, fill=codebackground, thick,
                minimum width=3.4cm, minimum height=1.25cm, text width=3.3cm, align=center, font=\scriptsize, inner sep=6pt},
    qbox/.style={rectangle, rounded corners=4pt, draw=navy!90!black, fill=navy!12, thick,
                 minimum width=3.0cm, minimum height=1.05cm, text width=2.9cm, align=center, font=\scriptsize\bfseries, inner sep=4pt},
    db/.style={cylinder, shape border rotate=90, aspect=0.22, draw=navy!80!black, fill=codebackground, thick,
               minimum width=3.4cm, minimum height=1.25cm, text width=3.3cm, inner sep=4pt, align=center, font=\scriptsize, text=navy!90!black},
    serpbox/.style={rectangle, rounded corners=4pt, draw=darkgreen!90!black, fill=darkgreen!12, thick,
                    minimum width=3.0cm, minimum height=1.25cm, align=center, font=\scriptsize\bfseries, inner sep=5pt},
    evaleff/.style={rectangle, rounded corners=4pt, draw=crimson!85!black, fill=crimson!8, thick,
                    font=\tiny\bfseries, text=crimson!90!black, align=center, inner sep=5pt},
    evalqual/.style={rectangle, rounded corners=4pt, draw=darkgreen!85!black, fill=darkgreen!8, thick,
                     font=\tiny\bfseries, text=darkgreen!90!black, align=center, inner sep=5pt},
    arr/.style={-Stealth, thick, navy!80!black, shorten >= 3pt, shorten <= 3pt}
]
    % Linea divisoria orizzontale nello spazio libero (y = 3.7)
    \draw[navy!50!black, dashed, thick] (-7.6, 3.7) -- (16.6, 3.7);
    \node[font=\bfseries\footnotesize, text=navy, anchor=south east] at (-2.4, 3.85) {PIPELINE ONLINE (Latenza \& Qualit\`a)};
    \node[font=\bfseries\footnotesize, text=navy!70!black, anchor=north east] at (-2.4, 3.55) {PIPELINE OFFLINE (Throughput \& I/O)};

    % Colonna 1: Query & Indexing (x = 0.0)
    \node[qbox] (query) at (0.0, 8.1) {Query Utente $q$};
    \node[qbox] (expquery) at (0.0, 6.6) {Expanded Query\\[2pt]{\tiny NLP / Intent}};
    \node[box] (qproc) at (0.0, 5.1) {\textbf{Query Processing}\\[3pt]{\tiny Candidate Selection}};
    \node[db] (invidx) at (0.0, 2.3) {\textbf{Inverted Index}\\[2pt]{\tiny Postings Compresse}};
    \node[box] (indexing) at (0.0, -0.8) {\textbf{Indexing Pipeline}\\[3pt]{\tiny Parsing e Tokenizzazione}};
    \node[db] (doccoll) at (0.0, -3.0) {\textbf{Document Collection}\\[2pt]{\tiny Web Crawl / Ingestion}};

    % Colonna 2: Features (x = 5.4)
    \node[box] (featlook) at (5.4, 5.1) {\textbf{Feature Lookup}\\[3pt]{\tiny Estrazione $\mathbf{x}_{q,d} \in \mathbb{R}^D$}};
    \node[db] (docfeat) at (5.4, 2.3) {\textbf{Document Features}\\[2pt]{\tiny PageRank, Segnali Statici}};
    \node[box] (featproc) at (5.4, -0.8) {\textbf{Feature Processor}\\[3pt]{\tiny Calcolo PageRank e Metadati}};

    % Colonna 3: Learning to Rank & SERP (x = 10.8)
    \node[box] (rankfunc) at (10.8, 5.1) {\textbf{Learned Ranker}\\[3pt]{\tiny Top-Score (GBRT/NN)}};
    \node[serpbox] (serp) at (14.8, 5.1) {Ranked SERP\\[2pt]{\tiny Top-$k$ Risultati}};
    \node[db] (ltrmodel) at (10.8, 2.3) {\textbf{Modello LtR}\\[2pt]{\tiny Parametri Addestrati}};
    \node[box] (training) at (10.8, -0.8) {\textbf{Training Offline}\\[3pt]{\tiny Ottimizzazione della Loss}};
    \node[db] (traindata) at (10.8, -3.0) {\textbf{Training Data}\\[2pt]{\tiny Query e Giudizi di Rilevanza}};

    % Connessioni Online (flusso orizzontale continuo e lookup verticali)
    \draw[arr] (query) -- (expquery);
    \draw[arr] (expquery) -- (qproc);
    \draw[arr] (invidx) -- (qproc);
    \draw[arr] (qproc) -- node[above, font=\tiny\itshape] {$\sim 10^3$ doc} (featlook);
    \draw[arr] (docfeat) -- (featlook);
    \draw[arr] (featlook) -- node[above, font=\tiny\itshape] {$\mathbf{x}_{q,d}$} (rankfunc);
    \draw[arr] (ltrmodel) -- (rankfunc);
    \draw[arr] (rankfunc) -- node[above, font=\tiny\itshape] {Top-$k$} (serp);

    % Connessioni Offline (ampio spazio e percorsi ariosi)
    \draw[arr] (doccoll) -- (indexing);
    \draw[arr] (indexing) -- (invidx);
    \draw[arr] (doccoll.east) -- (2.7, -3.0) -- (2.7, -0.8)
        node[midway, left=4pt, font=\tiny\itshape, text=navy!80!black] {Documenti grezzi}
        -- (featproc.west);
    \draw[arr] (featproc) -- node[pos=0.45, right=5pt, font=\tiny\itshape, text=navy!80!black] {PageRank / metadati} (docfeat);
    \draw[arr] (traindata) -- (training);
    \draw[arr] (training) -- (ltrmodel);

    % Badge Metriche: colonna dedicata a sinistra per efficienza, a destra per efficacia
    \node[evaleff] (ev_lat) at (-4.6, 5.1) {\textbf{Valutazione Efficienza:}\\[2pt]Latenza Query ($p_{95}, p_{99} < 50\text{ ms}$)\\[1pt]Throughput (QPS)};
    \draw[crimson, thick, dashed] (ev_lat) -- (qproc);

    \node[evaleff] (ev_idx) at (-4.6, -0.8) {\textbf{Valutazione Efficienza:}\\[2pt]Throughput Indicizzazione\\[1pt]($\text{doc}/\text{h}$, $\text{MB}/\text{s}$)};
    \draw[crimson, thick, dashed] (ev_idx) -- (indexing);

    \node[evalqual] (ev_eff) at (14.8, 2.3) {\textbf{Valutazione Efficacia:}\\[2pt]Precision@$k$, MAP, NDCG\\[1pt]User Happiness / CTR};
    \draw[darkgreen, thick, dashed] (serp) -- (ev_eff);
\end{tikzpicture}%
}
\caption{Architettura integrata di un moderno motore di ricerca e mappatura dei punti di rilevazione delle metriche di valutazione: efficienza offline (throughput di indicizzazione), efficienza online (latenza di query processing) ed efficacia qualitativa sulla SERP (NDCG, MAP, soddisfazione dell'utente).}
\label{fig:architettura-valutazione}
\end{figure}

\FloatBarrier
\subsection{Misurare la Soddisfazione dell'Utente (\textit{User Happiness})}
\label{subsec:misurare-felicita-utente}

Nel mondo operativo dei motori di ricerca commerciali, l'obiettivo finale consiste nel massimizzare la soddisfazione soggettiva dell'utente (\textit{User Happiness}). Poiché lo stato di appagamento cognitivo non \`e direttamente osservabile dall'infrastruttura hardware, i sistemi reali monitorano una combinazione di \textbf{proxy comportamentali online}:

\begin{itemize}
    \item \textbf{Clic e Click-Through Rate (CTR):} Proporzione di volte in cui un risultato presentato sulla SERP viene cliccato rispetto al numero totale di visualizzazioni.
    \begin{itemize}
        \item \textit{Criticit\`a dei titoli ingannevoli:} Titoli accattivanti o sintesi fuorvianti (\textit{clickbait}) generano un CTR artificialmente elevato, benché la pagina di destinazione si riveli completamente inutile per l'utente.
        \item \textit{La ricerca a zero clic (\textit{Zero-Click Search}):} L'assenza di clic non costituisce sistematicamente un indice di fallimento. Quando la query esprime un bisogno informativo fattuale (ad esempio il meteo, una definizione dizionariale, il risultato di una partita o una conversione di valuta), la visualizzazione immediata della risposta corretta tramite \textit{Knowledge Card} o \textit{Featured Snippet} soddisfa perfettamente l'utente senza richiedere alcuna navigazione esterna.
    \end{itemize}
    \item \textbf{Acquisti e Conversioni Economiche:} Nei motori per l'e-commerce, la metrica suprema risiede nella quota di ricerche che sfociano in un acquisto effettivo e nel volume economico transato dall'utente.
    \item \textbf{Fidelizzazione e Visite Ripetute:} Frequenza con cui gli utenti ritornano a impiegare la piattaforma nell'arco di una settimana o di un mese.
    \item \textbf{Dwell Time (Tempo di Permanenza):} Intervallo temporale che intercorre tra il clic su un risultato e il ritorno dell'utente alla SERP (\textit{bouncing back}). Un Dwell Time prolungato segnala che la risorsa era ricca di contenuti utili, mentre un rimbalzo immediato sulla SERP testimonia che il documento era irrilevante o inadeguato.
\end{itemize}

\FloatBarrier
\subsection{La Valutazione nell'Era dell'Agentic Search e dei Modelli Generativi}
\label{subsec:valutazione-era-generativa}

\begin{alertblock}{Riflessione Metodologica: La Felicit\`a dell'Utente con i Sistemi Generativi}
L'affermazione di interfacce conversazionali avanzate basate su Large Language Model (quali ChatGPT o motori ibridi Agentic Search) mette in crisi l'impianto tradizionale fondato su CTR e Dwell Time: l'utente non scorre pi\`u una lista di link esterni, ma interagisce con un agente in grado di sintetizzare risposte complete e orchestrare chiamate a tool esterni.

In tale orizzonte emergono nuove sfide e metriche di valutazione:
\begin{itemize}
    \item \textbf{Numero di turni conversazionali come proxy di insoddisfazione:} In un contesto dialogico, una sessione prolungata con molteplici iterazioni e raffinamenti non indica necessariamente un elevato gradimento, ma pu\`o denotare l'incapacit\`a del sistema di cogliere l'intento iniziale dell'utente, costringendolo a ripetuti chiarimenti correttivi.
    \item \textbf{Sensibilit\`a al prompting e varianza statistica:} Mentre una query IR lessicale produce un comportamento deterministico sull'indice, lievi variazioni nella formulazione di un prompt o fluttuazioni probabilistiche della temperatura di campionamento possono stravolgere la completezza e la veridicit\`a dell'output generativo.
    \item \textbf{Dicotomia tra reperimento e sintesi (Il Paradigma RAG):} Nei sistemi RAG (\textit{Retrieval-Augmented Generation}), la valutazione deve scorporare la qualit\`a del recupero documentale (accuratezza del retrieval top-$k$) dalla fedelt\`a della sintesi testuale generata (assenza di allucinazioni e rigorosa attribuzione alle fonti citate).
\end{itemize}
\end{alertblock}

\FloatBarrier
\section{Valutazione della Rilevanza e Metodologie di Annotazione}
\label{sec:rilevanza-annotazione}

Poiché i segnali online (clic, conversioni) risentono pesantemente del \textit{Position Bias} (gli utenti tendono a cliccare i primi risultati a prescindere dal loro effettivo valore) e non sono riproducibili sperimentalmente prima del rilascio in produzione di un nuovo modello, la misura regina per confrontare algoritmi di Information Retrieval rimane la \textbf{rilevanza intrinseca} dei documenti reperiti rispetto a un'esigenza informativa data~\cite{manning2008}.

\FloatBarrier
\subsection{La Rilevanza come Proxy Fondamentale e gli Esperimenti di Cranfield}
\label{subsec:cleverdon-cranfield}

L'idea di fondare la valutazione dell'IR su un banco di prova controllato fu introdotta da Cyril Cleverdon nei celebri \textbf{esperimenti di Cranfield} (Inghilterra, fine anni '50 e anni '60)~\cite{cleverdon1967}. 
In quell'epoca pionieristica non esistevano motori Web né calcolatori elettronici interattivi: Cleverdon operava su collezioni cartacee specializzate in aeronautica, confrontando l'efficacia di diversi linguaggi di indicizzazione e schede perforate attraverso interrogazioni reali formulate da ingegneri.

Il contributo concettuale di Cleverdon si \`e consolidato nella nozione di \textbf{Collezione di Test} (\textit{Test Collection}):

\begin{definition}[Collezione di Test Benchmark]
Una collezione di test standard per l'Information Retrieval \`e definita formalmente dalla tripletta:
\begin{equation}
\mathcal{T} = (\mathcal{D}, \mathcal{Q}, \mathcal{J})
\end{equation}
dove:
\begin{enumerate}
    \item $\mathcal{D} = \{d_1, d_2, \dots, d_N\}$ \`e la collezione finita di documenti benchmark.
    \item $\mathcal{Q} = \{q_1, q_2, \dots, q_M\}$ \`e una suite rappresentativa di bisogni informativi e query di prova.
    \item $\mathcal{J} \colon \mathcal{Q} \times \mathcal{D} \to \mathcal{S}$ \`e l'insieme dei giudizi di rilevanza (\textit{relevance judgments} o \textit{qrels}), che associa a ciascuna coppia query-documento un valore di pertinenza appartenente a un insieme discreto $\mathcal{S}$ (binario $\mathcal{S}=\{0, 1\}$ o graduato $\mathcal{S}=\{0, 1, \dots, m\}$).
\end{enumerate}
\end{definition}

\FloatBarrier
\subsection[Bisogno Informativo vs. Query]{Bisogno Informativo\\
(\textit{Information Need}) vs. Query}
\label{subsec:information-need-vs-query}

Un principio metodologico assoluto della teoria dell'IR risiede nella netta demarcazione tra l'intento mentale dell'essere umano e la sequenza di simboli immessa nella macchina:

\begin{itemize}
    \item \textbf{Bisogno Informativo (\textit{Information Need}):} \`E lo stato cognitivo dell'utente, l'esigenza o la lacuna di conoscenza che egli intende colmare. Esempio reale: \textit{``Il fondo della mia piscina si sta macchiando di nero e desidero capire come pulirlo senza svuotarla''}.
    \item \textbf{Query:} \`E la stringa di testo sintetica e spesso imperfetta che l'utente compone nel campo di ricerca per estrinsecare il proprio bisogno. Esempio: \texttt{``pulitore piscina fondo''}.
\end{itemize}

\begin{property}[Principio di Valutazione della Rilevanza]
La rilevanza di un documento deve essere rigorosamente valutata rispetto all'\textbf{Information Need} originario dell'utente, e \textbf{mai} rispetto alla mera rispondenza letterale dei token presenti nella query. Un documento che ripete ossessivamente le parole chiave della query ma non risponde al quesito sottostante \`e da considerarsi irrilevante.
\end{property}

\FloatBarrier
\subsection{La Sfida della Scalabilit\`a Combinatoria}
\label{subsec:scalabilita-combinatoria}

Nelle collezioni Web e aziendali contemporanee, l'annotazione esaustiva della funzione $\mathcal{J}$ si scontra con una barriera combinatoria insormontabile. Si consideri uno scenario operativo standard:
\begin{equation}
|\mathcal{D}| = 5 \times 10^6 \text{ documenti}, \quad |\mathcal{Q}| = 50\,000 \text{ query}
\end{equation}
Il numero totale di coppie query-documento da esaminare risulta:
\begin{equation}
|\mathcal{D}| \times |\mathcal{Q}| = 5 \times 10^6 \times 50\,000 = 2.5 \times 10^{11} = 0.25 \times 10^{12} \text{ coppie}
\end{equation}
ovvero \textbf{un quarto di trilione di giudizi}!

Ipotizzando che un annotatore umano impieghi in media $2.5\text{ secondi}$ per leggere e qualificare la pertinenza di ciascun documento rispetto al topic:
\begin{align}
\text{Tempo Complessivo} &= 2.5 \times 10^{11} \times 2.5\text{ s} = 6.25 \times 10^{11}\text{ secondi} \nonumber \\
&\approx 173\,611\,111\text{ ore/persona} \nonumber \\
&\approx 20\,000\text{ anni/uomo continuativi!}
\end{align}
Questo vincolo fisico dimostra l'impossibilit\`a teorica e pratica di procedere a una valutazione esaustiva dell'intero archivio, rendendo obbligatorio restringere il giudizio a un sottoinsieme accuratamente campionato di documenti per ciascuna query.

\FloatBarrier
\subsection{Provenienza e Rappresentativit\`a delle Query di Test}
\label{subsec:provenienza-query-test}

La scelta della suite $\mathcal{Q}$ deve garantire che le interrogazioni rispecchino la reale distribuzione delle esigenze informative degli utenti finali~\cite{manning2008}:

\begin{itemize}
    \item \textbf{Campionamento da Query Log Reali (Motori Web Moderni):} Sistemi commerciali quali Google, Microsoft Bing e Yahoo campionano le query di test direttamente dai propri log di transazione storici (come l'AOL Query Log o collezioni Microsoft ORCAS). Tale strategia preserva la validit\`a ecologica del test, catturando sia le query generiche ad altissimo traffico (\textit{head queries}) sia le ricerche complesse della coda lunga (\textit{tail queries}).
    \item \textbf{Definizione Manuale da Esperti di Dominio (Approccio Classico):} Nelle collezioni in cui non si dispone di registrazioni d'uso massive (o in contesti iperspecialistici), specialisti del settore redigono a tavolino scenari strutturati d'uso (\textit{topics}) e le correlate espressioni di ricerca.
    \item \textbf{Fallacia del Campionamento Sintetico:} Generare query campionando casualmente sequenze di termini o frasi dai documenti della collezione costituisce un errore metodologico gravissimo: le stringhe risultanti riflettono unicamente il lessico dei testi archiviati e non l'intento di ricerca umano, distorcendo irreparabilmente la valutazione del ranking.
\end{itemize}

\FloatBarrier
\subsection{Metodologie di Annotazione: Crowdsourcing ed Esperti}
\label{subsec:metodologie-annotazione}

La compilazione dei giudizi di rilevanza può seguire diverse strategie operative:

\begin{itemize}
    \item \textbf{Annotatori Esperti Dedicati:} Professionisti formati specificamente sul dominio di applicazione. Garantiscono la massima coerenza e precisione, ma comportano costi elevatissimi e bassa scalabilit\`a.
    \item \textbf{Piattaforme di Crowdsourcing (es. Amazon Mechanical Turk):} Sottomettono coppie query-documento a una moltitudine di lavoratori distribuiti a basso costo. Sebbene abbatta i costi monetari, la letteratura evidenzia una \textbf{varianza statistica molto elevata}, dovuta alla presenza di utenti frettolosi, incompetenti o spammer. Per ottenere dati stabili occorre sottoporre ogni coppia a molteplici valutatori indipendenti e applicare metriche di accordo inter-annotatore (quali la $\kappa$ di Cohen o la $\kappa$ di Fleiss) e filtri di maggioranza.
    \item \textbf{Valutazione Tramite LLM (\textit{LLM-as-a-Judge}):} L'impiego di modelli linguistici di frontiera per annotare la pertinenza query-documento sta rivoluzionando la velocit\`a di generazione dei benchmark. Tuttavia, occorre monitorare e correggere i noti \textit{bias intrinseci} degli LLM: il \textit{length bias} (tendenza a premiare testi verbosi e lunghi a prescindere dal contenuto), il \textit{self-preference bias} (preferenza per risposte stilisticamente affini a quelle del modello stesso) e il \textit{position bias} nel prompt.
\end{itemize}

\FloatBarrier
\section{Evoluzione Storica e Iniziative Attuali: Il Framework TREC}
\label{sec:storia-trec-benchmark}

\FloatBarrier
\subsection{Dalle Collezioni Pioniere all'Esplosione Dimensionale}
\label{subsec:corpora-storici}

La Tabella~\ref{tab:test-corpora-it} documenta l'evoluzione storica dei principali benchmark di valutazione in Information Retrieval, evidenziando il drastico passaggio di scala avvenuto negli anni '90.

\begin{table}[H]
\centering
\small
\resizebox{\columnwidth}{!}{%
\begin{tabular}{l r r r r r}
\toprule
\textbf{Collezione} & \textbf{Documenti ($N$)} & \textbf{Query ($M$)} & \textbf{Dimensione} & \textbf{Termini/Doc} & \textbf{Valutazioni (Q-D)} \\
\midrule
ADI & 82 & 35 & $< 1\text{ MB}$ & --- & Esaustive \\
AIT & 2\,109 & 14 & 2 MB & 400.0 & $> 10\,000$ \\
CACM & 3\,204 & 64 & 2 MB & 24.5 & Parziali \\
CISI & 1\,460 & 112 & 2 MB & 46.5 & Parziali \\
\textbf{Cranfield} & \textbf{1\,400} & \textbf{225} & \textbf{2 MB} & \textbf{53.1} & \textbf{Esaustive ($100\%$)} \\
LISA & 5\,872 & 35 & 3 MB & --- & Parziali \\
Medline & 1\,033 & 30 & 1 MB & --- & Esaustive \\
NPL & 11\,429 & 93 & 3 MB & --- & Parziali \\
OSHMED & 348\,566 & 106 & 400 MB & 250.0 & 16\,140 \\
Reuters-21578 & 21\,578 & 672 & 28 MB & 131.0 & Parziali \\
\midrule
\textbf{TREC (Tipica)} & \textbf{740\,000} & \textbf{200} & \textbf{2\,000 MB (2 GB)} & \textbf{89--3\,543} & \textbf{$> 100\,000$ (Pool)} \\
\bottomrule
\end{tabular}%
}
\caption{Panoramica dimensionale dei corpora benchmark storici e raffronto con l'architettura TREC~\cite{manning2008}.}
\label{tab:test-corpora-it}
\end{table}

Nelle collezioni pioniere degli anni '60 e '70 (come Cranfield e Medline), le ridotte dimensioni (poche migliaia di articoli scientifici e meno di $2\text{ MB}$ di dati) permettevano agli annotatori di analizzare fisicamente ogni singolo documento della collezione per ciascuna query, compilando una matrice di rilevanza completa al $100\%$. Con l'avvento di archivi da centinaia di migliaia di documenti e gigabyte di testo, questo approccio \`e divenuto impraticabile, determinando la nascita della Text REtrieval Conference.

\FloatBarrier
\subsection{L'Infrastruttura TREC (\textit{Text REtrieval Conference})}
\label{subsec:infrastruttura-trec}

Istituita nel 1992 e sponsorizzata congiuntamente dal National Institute of Standards and Technology (NIST) e dal Dipartimento della Difesa statunitense (DARPA), la conferenza \textbf{TREC} ha fornito alla comunit\`a scientifica internazionale l'infrastruttura metodologica e computazionale per testare sistemi IR su grandissima scala~\cite{manning2008}.

L'organizzazione di TREC si fonda su tre pilastri:
\begin{enumerate}
    \item \textbf{Tracce Tematiche (\textit{Tracks}):} Aree di ricerca focalizzate su specifici compiti operativi o domini applicativi (ricerca Web, e-discovery legale, genomica e biomedicina, question answering, deep retrieval, enterprise search).
    \item \textbf{Definizione Formale dei Topic:} A differenza delle brevi query del Web, un \textbf{Topic TREC} descrive formalmente l'esigenza informativa mediante tre campi distinti:
    \begin{itemize}
        \item \textbf{Title:} Una brevissima sequenza di parole chiave (simile a una query Web reale, es. \texttt{``acquisizione idrocarburi artico''}).
        \item \textbf{Description:} Una frase che sintetizza l'obiettivo del bisogno informativo (es. \textit{``Identificare accordi commerciali per lo sfruttamento di giacimenti di petrolio o gas nel Mare Glaciale Artico''}).
        \item \textbf{Narrative:} Un testo esplicativo redatto dal committente che specifica in dettaglio cosa renda un documento rilevante e cosa debba essere rigorosamente escluso (es. \textit{``I documenti pertinenti devono citare trattati o contratti stipulati tra governi o corporazioni. Sono esclusi studi meramente geologici o ambientali privi di risvolti economici diretti''}).
    \end{itemize}
    \item \textbf{Il Concetto di Run:} Il file contenente i risultati (tipicamente i primi $1\,000$ documenti per ciascun topic, ordinati per punteggio decrescente) inviato da ciascun gruppo di ricerca prima della scadenza della gara prende il nome di \textbf{Run}.
\end{enumerate}

\FloatBarrier
\subsection{La Tecnica del Pooling e l'Ipotesi dei Documenti Non Giudicati}
\label{subsec:tecnica-pooling}

Per risolvere l'impasse combinatoria senza rinunciare all'affidabilit\`a statistica, le prime edizioni di TREC hanno introdotto la tecnica del \textbf{Pooling} (Figura~\ref{fig:pooling-process}):

\begin{figure}[H]
\centering
\resizebox{0.95\linewidth}{!}{%
\begin{tikzpicture}[
    box/.style={rectangle, rounded corners=4pt, draw=navy, fill=codebackground, thick, 
                minimum width=2.4cm, minimum height=0.85cm, text width=2.3cm, align=center, font=\small},
    pool/.style={rectangle, rounded corners=5pt, draw=navy!90!black, fill=navy!10, thick, 
                 minimum width=3.6cm, minimum height=2.3cm, text width=3.4cm, align=center, font=\small, inner sep=5pt},
    eval/.style={rectangle, rounded corners=5pt, draw=darkgreen, fill=darkgreen!10, thick, 
                 minimum width=3.2cm, minimum height=2.3cm, text width=3.1cm, align=center, font=\small, inner sep=5pt}
]
    \node[box] (sys1) at (-5.6, 1.5)  {Sistema $S_1$\\\footnotesize (Run 1)};
    \node[box] (sys2) at (-5.6, 0.0)  {Sistema $S_2$\\\footnotesize (Run 2)};
    \node[box] (sysk) at (-5.6, -1.5) {Sistema $S_m$\\\footnotesize (Run $m$)};

    \node[pool] (poolbox) at (-0.7, 0.0) {
        \textbf{Pool di Candidati $\mathcal{P}$}\\[3pt]
        $\mathcal{P} = \bigcup_{i=1}^m \text{Top-}k(S_i)$\\[4pt]
        \footnotesize\textcolor{navy!80!black}{(unione priva di duplicati)}
    };

    \node[eval] (assess) at (5.2, 0.0) {
        \textbf{Annotatori Umani}\\[4pt]
        Valutazione manuale\\
        della rilevanza\\[4pt]
        \footnotesize\textit{(Relevant / Non-Relevant)}
    };

    \draw[-Stealth, thick, navy] (sys1.east) -- node[above, sloped, font=\scriptsize] {Top-$k$} (poolbox.155);
    \draw[-Stealth, thick, navy] (sys2.east) -- node[above, font=\scriptsize] {Top-$k$} (poolbox.west);
    \draw[-Stealth, thick, navy] (sysk.east) -- node[below, sloped, font=\scriptsize] {Top-$k$} (poolbox.205);
    \draw[-Stealth, very thick, darkgreen] (poolbox.east) -- node[above=2pt, font=\scriptsize\bfseries, text=darkgreen, align=center] {Solo documenti\\di $\mathcal{P}$} (assess.west);
\end{tikzpicture}%
}
\caption{Architettura operativa del processo di Pooling: i primi $k$ documenti estratti da ciascuna run indipendente confluiscono nel pool deduplicato $\mathcal{P}$ sottoposto ad annotazione umana.}
\label{fig:pooling-process}
\end{figure}

Il workflow operativo del pooling procede come segue:
\begin{enumerate}
    \item Si ricevono le run prodotte da un insieme variegato di $m$ sistemi partecipanti (costruiti con paradigmi algoritmici eterogenei: booleano esteso, BM25, vettoriale VSM, modelli probabilistici e di espansione query).
    \item Per ogni query, si prelevano i primi $k$ documenti restituiti da ciascun sistema (tipicamente $k = 50$ o $k = 100$).
    \item Si calcola l'unione insiemistica di questi elenchi, rimuovendo tutti i duplicati per costituire il pool dei candidati:
    \begin{equation}
    \mathcal{P}(q) = \bigcup_{i=1}^m \text{Top-}k(S_i, q)
    \end{equation}
    \item Gli annotatori umani del NIST esaminano e giudicano la pertinenza \textbf{soltanto} dei documenti compresi nel pool $\mathcal{P}(q)$.
\end{enumerate}

\begin{alertblock}{L'Assunzione dei Documenti Non Giudicati (\textit{Unpooled Documents})}
Tutti i documenti dell'archivio $\mathcal{D}$ che non compaiono nel pool $\mathcal{P}(q)$ vengono considerati convenzionalmente \textbf{non rilevanti} ($r = 0$).

\textbf{Validit\`a empirica e limiti:}
Sebbene questa assunzione introduca un potenziale \textit{pooling bias} a sfavore di algoritmi innovativi futuri che dovessero recuperare documenti rilevanti non intercettati da alcuno dei sistemi originari, le verifiche sperimentali condotte da NIST hanno dimostrato che:
\begin{itemize}
    \item La stragrande maggioranza dei documenti pertinenti della collezione viene intercettata purché il pool includa run generate da modelli sufficientemente diversificati.
    \item Sebbene le metriche assolute di un nuovo sistema possano risultare leggermente sottostimate, il \textbf{ranking comparativo} tra sistemi rimane eccezionalmente stabile ed equo.
\end{itemize}
\end{alertblock}

\FloatBarrier
\section{Metriche di Valutazione Binarie per Insiemi Non Ordinati}
\label{sec:metriche-binarie-non-ordinate}

Nel paradigma dell'\textit{Unranked Retrieval} (tipico dei sistemi booleani esatti), il motore restituisce un insieme piatto di documenti privo di ordinamento qualitativo. Ciascun documento della collezione \`e bipartito rispetto alla rilevanza effettiva e allo stato di recupero.

\FloatBarrier
\subsection{Matrice di Contingenza e Spazio delle Decisioni}
\label{subsec:matrice-contingenza}

La valutazione si fonda sulla scomposizione dello spazio campionario nella classica \textbf{matrice di contingenza $2 \times 2$} (Tabella~\ref{tab:contingency-evaluation}):

\begin{table}[H]
\centering
\small
\begin{tabularx}{0.88\columnwidth}{l >{\centering\arraybackslash}X >{\centering\arraybackslash}X}
\toprule
& \textbf{Documento Rilevante} & \textbf{Documento Non Rilevante} \\
\midrule
\textbf{Documento Reperito} & True Positive ($TP$) & False Positive ($FP$) \\
\addlinespace
\textbf{Documento Non Reperito} & False Negative ($FN$) & True Negative ($TN$) \\
\bottomrule
\end{tabularx}
\caption{Matrice di contingenza per la classificazione query-documento.}
\label{tab:contingency-evaluation}
\end{table}

La Figura~\ref{fig:venn-retrieval} offre una chiara lettura insiemistica mediante diagramma di Eulero-Venn.

\begin{figure}[H]
\centering
\resizebox{0.85\columnwidth}{!}{%
\begin{tikzpicture}
    % Universo D
    \draw[thick, fill=codebackground, draw=navy!40, rounded corners=8pt] (-6.4, -4.2) rectangle (6.4, 4.2);
    \node[anchor=north west, font=\bfseries\small, text=navy] at (-6.0, 3.9) {Universo $\boldsymbol{\mathcal{D}}$ \normalfont\footnotesize\textcolor{navy!75!black}{(Collezione Documentale)}};

    % Insiemi Rilevanti e Reperiti (centri a y = -0.35, raggio 2.2cm, sommita a y = 1.85)
    \draw[thick, fill=navy!15, draw=navy, fill opacity=0.6] (-1.45, -0.35) circle (2.2cm);
    \draw[thick, fill=darkgreen!18, draw=darkgreen, fill opacity=0.6] (1.45, -0.35) circle (2.2cm);

    % Titoli insiemi (ancorati a south a y = 2.15: base a 3mm sopra la sommita dei cerchi e oltre 7mm sotto il titolo dell'Universo)
    \node[anchor=south, font=\bfseries\normalsize, text=navy, align=center] at (-1.8, 2.15) {Rilevanti ($\mathcal{R}$)\\[2pt]\scriptsize\textit{(Pertinenti al bisogno)}};
    \node[anchor=south, font=\bfseries\normalsize, text=darkgreen, align=center] at (1.8, 2.15) {Reperiti ($\mathcal{A}$)\\[2pt]\scriptsize\textit{(Restituiti dal motore)}};

    % Regioni interne (solo sigla metrica e formula insiemistica pulita)
    \node[align=center, text=navy!90!black] at (-2.5, -0.35) {\textbf{\Large FN}\\[4pt]\normalsize $\mathcal{R} \setminus \mathcal{A}$};
    \node[align=center, text=navy!90!black] at (0, -0.35) {\textbf{\Large TP}\\[4pt]\normalsize $\mathcal{R} \cap \mathcal{A}$};
    \node[align=center, text=darkgreen!90!black] at (2.5, -0.35) {\textbf{\Large FP}\\[4pt]\normalsize $\mathcal{A} \setminus \mathcal{R}$};
    
    % Legenda descrittiva delle regioni (disposta a y = -3.05, 5mm sotto i cerchi)
    \node[font=\scriptsize, text=navy!90!black, align=center] at (-3.7, -3.05) {\textbf{FN}: Rilevanti non reperiti};
    \node[font=\scriptsize, text=black, align=center] at (0, -3.05) {\textbf{TP}: Reperiti rilevanti};
    \node[font=\scriptsize, text=darkgreen!90!black, align=center] at (3.7, -3.05) {\textbf{FP}: Reperiti non rilevanti};

    % True Negative (disposto a y = -3.75, perfettamente allineato sulla riga base)
    \node[font=\scriptsize, text=gray!80!black, align=center] at (0, -3.75) {\textbf{TN (True Negative)}: Documenti non rilevanti e non reperiti dal motore \quad $(\mathcal{D} \setminus (\mathcal{R} \cup \mathcal{A}))$};
\end{tikzpicture}%
}
\caption{Rappresentazione insiemistica del reperimento documentale e partizione delle regioni della matrice di confusione.}
\label{fig:venn-retrieval}
\end{figure}

\FloatBarrier
\subsection{Precisione, Richiamo e loro Trade-off}
\label{subsec:precision-recall-tradeoff}

\begin{definition}[Precisione (\textit{Precision}, $P$)]
Frazione dei documenti restituiti dal motore che risultano effettivamente rilevanti rispetto al bisogno informativo:
\begin{equation}
P = \frac{TP}{TP + FP} = P(\text{rilevante} \mid \text{reperito})
\end{equation}
\end{definition}

\begin{definition}[Richiamo (\textit{Recall}, $R$)]
Frazione dei documenti rilevanti complessivamente presenti nella collezione che il sistema \`e stato in grado di individuare ed estrarre:
\begin{equation}
R = \frac{TP}{TP + FN} = P(\text{reperito} \mid \text{rilevante})
\end{equation}
\end{definition}

Tra Precisione e Richiamo sussiste un \textbf{trade-off intrinseco}:
\begin{itemize}
    \item Un sistema banale che restituisce l'intero archivio documentale ($\mathcal{A} = \mathcal{D}$) raggiunge un richiamo perfetto $R = 1.0$ ($100\%$), ma la sua precisione collassa a valori irrisori ($P \approx 0$).
    \item All'opposto, un sistema iper-selettivo che restituisce esclusivamente il singolo documento in assoluto pi\`u probabile può ottenere $P = 1.0$ ($100\%$), sacrificando quasi integralmente il richiamo ($R \approx 0$).
\end{itemize}

\FloatBarrier
\subsection{La Media Armonica: F-Measure ($F_1$-Score)}
\label{subsec:f-measure-armonica}

Per sintetizzare la bont\`a di un sistema in un indicatore scalare unico bilanciando $P$ ed $R$, si definisce la \textbf{F-Measure} (o $F_1$-score) come la \textbf{media armonica} di Precisione e Richiamo:
\begin{equation}
F_1 = \frac{2}{\frac{1}{P} + \frac{1}{R}} = \frac{2 \cdot P \cdot R}{P + R} = \frac{2 \cdot TP}{2 \cdot TP + FP + FN}
\end{equation}

\begin{property}[Propriet\`a della Media Armonica e Superiorit\`a rispetto alla Media Aritmetica]
La media armonica si rivela l'unica scelta rigorosa per la valutazione congiunta in virtù delle sue propriet\`a:
\begin{enumerate}
    \item \`E sempre dominata dalla media geometrica e dalla media aritmetica:
    \begin{equation}
    \frac{2}{\frac{1}{P} + \frac{1}{R}} \le \sqrt{P \cdot R} \le \frac{P + R}{2}
    \end{equation}
    \item Quando i due parametri divergono in modo marcato, la media armonica collassa verso il \textbf{minimo} tra i due valori, punendo severamente ogni squilibrio algoritmico.
\end{enumerate}
\end{property}

\begin{example}[La Fallacia della Media Aritmetica]
Si consideri il motore banale che restituisce l'intera collezione: $R = 1.0$ e $P = 0.0001$.
\begin{itemize}
    \item La media aritmetica assegna un punteggio assurdo superiore al $50\%$:
    \begin{equation}
    \text{Media Aritmetica} = \frac{1.0 + 0.0001}{2} \approx 0.50005 \quad (50.0\%)
    \end{equation}
    permettendo a un sistema totalmente inservibile di apparire mediocremente efficace.
    \item La media armonica ($F_1$) demolisce immediatamente il sistema:
    \begin{equation}
    F_1 = \frac{2 \cdot 1.0 \cdot 0.0001}{1.0 + 0.0001} \approx 0.0002 \quad (0.02\%)
    \end{equation}
    restituendo una misura realistica dell'inutilit\`a del recupero.
\end{itemize}
\end{example}

\FloatBarrier
\subsection{La Misura Pesata $F_\beta$}
\label{subsec:f-beta-pesata}

Se l'applicazione impone di dare priorit\`a asimmetrica a una delle due grandezze, si introduce un fattore di ponderazione $\alpha \in [0, 1]$:
\begin{equation}
F = \frac{1}{\alpha \frac{1}{P} + (1 - \alpha) \frac{1}{R}}
\end{equation}
Ponendo convenientemente $\alpha = \frac{1}{1 + \beta^2}$, si ricava la formulazione canonica di \textbf{$F_\beta$}:
\begin{equation}
F_\beta = \frac{(1 + \beta^2) \cdot P \cdot R}{\beta^2 \cdot P + R}
\end{equation}

Il parametro $\beta \in [0, +\infty)$ controlla il bilanciamento relativo:
\begin{itemize}
    \item $\beta = 1.0$: Si riottiene la classica $F_1$ a pesatura identica.
    \item $\beta < 1.0$ (ad esempio $\beta = 0.5$): Il modello valorizza la \textbf{Precisione} rispetto al Richiamo. Tipico della ricerca Web generica, dove all'utente interessa trovare immediatamente pochi documenti perfetti.
    \item $\beta > 1.0$ (ad esempio $\beta = 2.0$): Il modello valorizza il \textbf{Richiamo} rispetto alla Precisione. Imprescindibile nelle ricerche brevettuali (\textit{patent search}), nelle perizie medico-legali e nell'e-discovery giuridico, dove mancare anche un solo documento rilevante comporta danni irreparabili.
\end{itemize}

\FloatBarrier
\section{Metriche Basate sul Ranking (\textit{Ranked Measures} - Rilevanza Binaria)}
\label{sec:metriche-ranking-binarie}

Nei motori di ricerca operativi i documenti non formano insiemi caotici, ma vengono presentati in una graduatoria ordinata in base al punteggio decrescente di rilevanza stimata. Poiché gli utenti esaminano quasi esclusivamente i primissimi risultati della SERP, la valutazione deve premiare la concentrazione dei documenti pertinenti in cima alla lista.

\FloatBarrier
\subsection{Precision@k e Recall@k: Valutazione alla Soglia di Rango}
\label{subsec:precision-at-k}

\begin{definition}[Precision@k ($P@k$)]
Data una soglia di rango $k \in \mathbb{N}^+$, $P@k$ corrisponde alla proporzione di documenti rilevanti presenti nei primi $k$ posti della graduatoria restituita:
\begin{equation}
P@k = \frac{1}{k} \sum_{i=1}^k \mathrm{rel}_i
\end{equation}
dove $\mathrm{rel}_i \in \{0, 1\}$ rappresenta l'esito binario di rilevanza del documento collocato al rango $i$.
\end{definition}

\begin{example}[Calcolo Discreto di $P@k$ ed Evidenza del Comportamento Non Monotono]
Sia data la graduatoria ordinata dei primi 5 documenti con rilevanze $\mathbf{r} = \langle 1, 1, 0, 0, 1 \rangle$:
\begin{itemize}
    \item $P@1 = \frac{1}{1} = 1.00$
    \item $P@2 = \frac{2}{2} = 1.00$
    \item $P@3 = \frac{2}{3} \approx 0.67$
    \item $P@4 = \frac{2}{4} = 0.50$ \quad (\textbf{si noti come $P@4 < P@3$})
    \item $P@5 = \frac{3}{5} = 0.60$ \quad (\textbf{la precisione torna a crescere})
\end{itemize}
\end{example}

In modo complementare, $Recall@k$ ($R@k$) calcola la frazione dei documenti rilevanti complessivi $R$ recuperati entro la soglia $k$:
\begin{equation}
R@k = \frac{1}{R} \sum_{i=1}^k \mathrm{rel}_i
\end{equation}

\FloatBarrier
\subsection{Limiti Metodologici di $P@k$ nella Media tra Query}
\label{subsec:limiti-precision-at-k}

\begin{alertblock}{La Distorsione di $P@k$ nella Media Macro}
Sebbene $P@k$ fornisca un'indicazione immediata della pulizia della prima schermata, presenta un difetto metodologico strutturale quando viene mediato su insiemi di query eterogenee:
\begin{itemize}
    \item Il numero totale $R$ di documenti rilevanti presenti nella collezione varia sensibilmente da query a query.
    \item Se per la query $q_A$ esistono soltanto $R_A = 2$ documenti rilevanti nell'intera collezione, il punteggio massimo teorico di $P@10(q_A)$ \`e pari a $\frac{2}{10} = 0.20$, persino qualora il motore sia infallibile e collochi entrambi i documenti esattamente ai primi due posti!
    \item Viceversa, per una query generica $q_B$ con $R_B = 100$, il sistema può conseguire comodamente $P@10(q_B) = 1.0$.
    \item Mediando aritmeticamente $P@10$, l'algoritmo viene penalizzato dalla scarsit\`a oggettiva dei documenti rilevanti di talune query e non da deficienze del proprio ordinamento.
\end{itemize}
\end{alertblock}

\FloatBarrier
\subsection{Average Precision (AP)}
\label{subsec:average-precision-ap}

L'\textbf{Average Precision} (AP) risolve i limiti delle metriche a soglia rigida, sintetizzando l'intera curva Precision-Recall per una singola query in un indicatore scalare unitario.

\begin{definition}[Average Precision]
Dato un ordinamento di $N$ documenti per una query avente $R$ documenti rilevanti totali presenti nell'archivio, l'AP \`e la media delle precisioni istantanee $P@k$ calcolate in corrispondenza di ciascuna posizione $k$ in cui si incontra un documento rilevante (ovvero dove il richiamo subisce un incremento):
\begin{equation}
AP = \frac{1}{R} \sum_{k=1}^N P@k \cdot \mathrm{rel}_k = \frac{1}{R} \sum_{k \in \{K_1, K_2, \dots, K_R\}} P@k
\end{equation}
dove $K_1, K_2, \dots, K_R$ indicano i ranghi dei documenti rilevanti reperiti.
\end{definition}

\begin{property}[Interpretazione Geometrica dell'Average Precision]
L'Average Precision possiede un solido fondamento matematico:
\begin{enumerate}
    \item In corrispondenza di ogni documento rilevante al rango $k$, il richiamo cresce di un incremento discreto pari a $\Delta R = \frac{1}{R}$. Pertanto, la sommatoria calcola la somma delle aree dei rettangoli $P@k \cdot \Delta R$, rappresentando l'esatta \textbf{Area Sotto la Curva Precision-Recall} (AUC-PR).
    \item La normalizzazione per $R$ (il totale assoluto dei documenti rilevanti nell'intera collezione, e non solo quelli estratti) impone che tutti i rilevanti non compresi nei primi $N$ ranghi concorrano implicitamente con precisione pari a $0$, penalizzando severamente i sistemi a basso richiamo.
\end{enumerate}
\end{property}

La Figura~\ref{fig:ap-blocks} illustra graficamente la sequenza di computazione dell'AP.

\begin{figure}[H]
\centering
\begin{tikzpicture}[
    scale=0.88,
    docrel/.style={rectangle, rounded corners=3pt, draw=darkgreen, fill=darkgreen!15, thick, minimum width=1.15cm, minimum height=0.65cm, font=\footnotesize\bfseries},
    docnon/.style={rectangle, rounded corners=3pt, draw=crimson, fill=crimson!12, thick, minimum width=1.15cm, minimum height=0.65cm, font=\footnotesize},
    lbl/.style={font=\scriptsize, align=center}
]
    \node[lbl] at (-1.6, 0.35) {Posizione ($k$):};
    \node[lbl] at (-1.6, -0.35) {Valutazione:};
    \node[lbl] at (-1.6, -1.05) {Precision@k:};

    \foreach \k/\rel/\p/\type in {
        1/1/1.00/docrel,
        2/0/---/docnon,
        3/1/0.67/docrel,
        4/0/---/docnon,
        5/1/0.60/docrel%
    } {
        \node[font=\footnotesize\bfseries] at (\k*1.65, 0.35) {$k=\k$};
        \node[\type] at (\k*1.65, -0.35) {\ifnum\rel=1 Rel\else NonRel\fi};
        \node[font=\footnotesize] at (\k*1.65, -1.05) {\ifnum\rel=1 \textbf{\p}\else \color{gray!60}--- \fi};
    }

    \draw[-Stealth, thick, navy] (0.4, 0.72) -- node[above, font=\scriptsize] {Ordinamento del Motore di Ricerca} (8.8, 0.72);
\end{tikzpicture}
\caption{Calcolo passo-passo di Average Precision per un ordinamento di 5 elementi con $R=3$ documenti rilevanti totali.}
\label{fig:ap-blocks}
\end{figure}

\begin{example}[Studio Comparativo tra Due Ordinamenti Alternativi]
Si consideri una query avente $R = 6$ documenti rilevanti nella collezione. Due motori alternativi producono i seguenti primi 10 documenti:
\begin{itemize}
    \item \textbf{Ranking \#1 (Rilevanti concentrati in cima):}
    \[
    \langle \text{Rel}, \text{NonRel}, \text{Rel}, \text{Rel}, \text{Rel}, \text{Rel}, \text{NonRel}, \text{NonRel}, \text{NonRel}, \text{Rel} \rangle
    \]
    \item \textbf{Ranking \#2 (Rilevanti ritardati e dispersi):}
    \[
    \langle \text{NonRel}, \text{Rel}, \text{NonRel}, \text{NonRel}, \text{Rel}, \text{Rel}, \text{Rel}, \text{NonRel}, \text{Rel}, \text{Rel} \rangle
    \]
\end{itemize}
La Tabella~\ref{tab:ap-comparison} traccia l'evoluzione di Precisione e Richiamo rango per rango.

\begin{table}[H]
\centering
\small
\resizebox{\columnwidth}{!}{%
\begin{tabular}{c c c c c c c}
\toprule
& \multicolumn{3}{c}{\textbf{Ranking \#1}} & \multicolumn{3}{c}{\textbf{Ranking \#2}} \\
\cmidrule(lr){2-4} \cmidrule(lr){5-7}
\textbf{Rango ($k$)} & \textbf{Esito} & \textbf{Recall} & \textbf{Precision} & \textbf{Esito} & \textbf{Recall} & \textbf{Precision} \\
\midrule
1  & \textbf{Rel}    & 0.17 & \textbf{1.00} & NonRel          & 0.00 & --- \\
2  & NonRel          & 0.17 & ---           & \textbf{Rel}    & 0.17 & \textbf{0.50} \\
3  & \textbf{Rel}    & 0.33 & \textbf{0.67} & NonRel          & 0.17 & --- \\
4  & \textbf{Rel}    & 0.50 & \textbf{0.75} & NonRel          & 0.17 & --- \\
5  & \textbf{Rel}    & 0.67 & \textbf{0.80} & \textbf{Rel}    & 0.33 & \textbf{0.40} \\
6  & \textbf{Rel}    & 0.83 & \textbf{0.83} & \textbf{Rel}    & 0.50 & \textbf{0.50} \\
7  & NonRel          & 0.83 & ---           & \textbf{Rel}    & 0.67 & \textbf{0.57} \\
8  & NonRel          & 0.83 & ---           & NonRel          & 0.67 & --- \\
9  & NonRel          & 0.83 & ---           & \textbf{Rel}    & 0.83 & \textbf{0.56} \\
10 & \textbf{Rel}    & 1.00 & \textbf{0.60} & \textbf{Rel}    & 1.00 & \textbf{0.60} \\
\midrule
\multicolumn{2}{l}{\textbf{Average Precision (AP)}} & & \textbf{0.78} & & & \textbf{0.52} \\
\bottomrule
\end{tabular}%
}
\caption{Confronto di Average Precision: Ranking \#1 premia la reattivit\`a posizionale anticipando i rilevanti.}
\label{tab:ap-comparison}
\end{table}

I rispettivi punteggi analitici valgono:
\begin{align}
AP(\text{Ranking \#1}) &= \frac{1.00 + 0.67 + 0.75 + 0.80 + 0.83 + 0.60}{6} = \frac{4.65}{6} \approx \mathbf{0.78} \\
AP(\text{Ranking \#2}) &= \frac{0.50 + 0.40 + 0.50 + 0.57 + 0.56 + 0.60}{6} = \frac{3.13}{6} \approx \mathbf{0.52}
\end{align}
Sebbene entrambi i motori recuperino tutti e 6 i documenti rilevanti entro il decimo rango ($R=100\%$), Ranking \#1 distanzia nettamente il concorrente ($0.78$ contro $0.52$), premiando il risparmio di sforzo visivo garantito all'utente.
\end{example}

\FloatBarrier
\subsection{Mean Average Precision (MAP)}
\label{subsec:mean-average-precision-map}

Quando si valuta l'efficacia di un algoritmo sull'intera suite di query benchmark $\mathcal{Q}$, si adotta la \textbf{Mean Average Precision} (MAP):

\begin{definition}[Mean Average Precision]
Media aritmetica dei valori di Average Precision calcolati su tutte le query del benchmark:
\begin{equation}
MAP = \frac{1}{|\mathcal{Q}|} \sum_{j=1}^{|\mathcal{Q}|} AP(q_j)
\end{equation}
\end{definition}

\begin{property}[Macro-Averaging e Tutela della Coda Lunga]
Il calcolo di MAP costituisce una procedura formale di \textbf{Macro-Averaging}: a ciascuna query viene assegnato esattamente lo stesso peso pari a $1/|\mathcal{Q}|$.
\begin{itemize}
    \item Se si impiegasse il \textit{Micro-Averaging} (sommare tutti i contributi di precisione dell'intero benchmark dividendo per il totale cumulativo dei documenti rilevanti esistenti), il punteggio globale verrebbe cannibalizzato da pochissime query ad altissima densit\`a documentale (es. \texttt{``sport''} o \texttt{``notizie''}, con decine di migliaia di rilevanti).
    \item Il Macro-Averaging preserva la parit\`a di rappresentanza: l'insuccesso su una query specialistica di coda lunga (\textit{tail query}) incide sul punteggio finale con la medesima gravit\`a del fallimento su una query comune.
\end{itemize}
\end{property}

\begin{example}[Calcolo di MAP su Due Query Distinte]
Si consideri un testbed con due bisogni informativi $q_1$ e $q_2$:
\begin{itemize}
    \item \textbf{Query $q_1$ ($R_1 = 5$ rilevanti totali):} Reperiti 5 rilevanti ai ranghi $\{1, 3, 6, 9, 10\}$ con precisioni istantanee $1.00$, $0.67$, $0.50$, $0.44$, $0.50$:
    \begin{equation}
    AP(q_1) = \frac{1.00 + 0.67 + 0.50 + 0.44 + 0.50}{5} = \frac{3.11}{5} \approx \mathbf{0.62}
    \end{equation}
    \item \textbf{Query $q_2$ ($R_2 = 3$ rilevanti totali):} Reperiti 3 rilevanti ai ranghi $\{2, 5, 7\}$ con precisioni istantanee $0.50$, $0.40$, $0.43$:
    \begin{equation}
    AP(q_2) = \frac{0.50 + 0.40 + 0.43}{3} = \frac{1.33}{3} \approx \mathbf{0.44}
    \end{equation}
\end{itemize}
Il valore di Mean Average Precision del sistema su questo benchmark \`e:
\begin{equation}
MAP = \frac{AP(q_1) + AP(q_2)}{2} = \frac{0.62 + 0.44}{2} = \mathbf{0.53}
\end{equation}
\end{example}

\FloatBarrier
\subsection{Metriche Troncate per il Web e Recommender Systems: AP@k e MAP@k}
\label{subsec:ap-map-at-k}

Nell'architettura del Web search moderno e nei sistemi di raccomandazione, calcolare l'AP sull'intero archivio \`e impraticabile: il numero totale $R$ di documenti rilevanti sparsi nel Web \`e ignoto, e all'utente finale non interessa ciò che accade oltre la ventesima posizione.

Si definiscono pertanto le versioni troncate a una soglia prestabilita $K$ (tipicamente $K=10$ o $K=20$):

\begin{definition}[AP@K e MAP@K]
Data la profondit\`a di taglio $K$, l'Average Precision at $K$ per la query $q_i$ \`e:
\begin{equation}
AP@K(q_i) = \frac{1}{K_i} \sum_{k=1}^K P@k(q_i) \cdot \mathrm{rel}(q_i, k)
\end{equation}
dove il fattore di normalizzazione $K_i$ corrisponde al numero effettivo di documenti rilevanti reperibili entro la soglia:
\begin{equation}
K_i = \min(R_i, K) \quad \left(\text{oppure } K_i = \sum_{k=1}^K \mathrm{rel}(q_i, k)\right)
\end{equation}
La metrica globale \textbf{MAP@K} su $\mathcal{Q}$ corrisponde alla media aritmetica:
\begin{equation}
MAP@K(\mathcal{Q}) = \frac{1}{|\mathcal{Q}|} \sum_{q_i \in \mathcal{Q}} AP@K(q_i)
\end{equation}
\end{definition}

Normalizzando per $\min(R_i, K)$ anziché per il valore rigido $K$, una query con soli 2 rilevanti totali che li collochi ai primi due posti consegue $AP@10 = 1.0$, eliminando la distorsione strutturale di $P@10$.

\FloatBarrier
\subsection{Mean Reciprocal Rank (MRR)}
\label{subsec:mean-reciprocal-rank-mrr}

In contesti applicativi quali:
\begin{itemize}
    \item \textit{Known-Item Search:} L'utente cerca una pagina o un file specifico di cui conosce l'esistenza;
    \item \textit{Query Navigazionali:} Ricerca dell'accesso o della homepage di un servizio (ad es. \texttt{``portale studenti unipi''});
    \item \textit{Factoid Question Answering:} Domande che ammettono una singola risposta esatta,
\end{itemize}
l'utente non ha alcun interesse a esaminare dieci documenti rilevanti, ma desidera localizzare unicamente la \textbf{prima risposta corretta}. Lo sforzo cognitivo e il tempo speso dall'utente sono inversamente proporzionali alla posizione di rango del primo risultato utile:

\begin{definition}[Reciprocal Rank e MRR]
Sia $\mathrm{rank}_1(q)$ la posizione ordinale del primo documento rilevante restituito per la query $q$. Il \textbf{Reciprocal Rank} (RR) \`e definito come:
\begin{equation}
RR(q) = \frac{1}{\mathrm{rank}_1(q)}
\end{equation}
Il \textbf{Mean Reciprocal Rank} (MRR) \`e la media aritmetica di $RR$ sull'insieme delle query benchmark:
\begin{equation}
MRR = \frac{1}{|\mathcal{Q}|} \sum_{i=1}^{|\mathcal{Q}|} \frac{1}{\mathrm{rank}_1(q_i)}
\end{equation}
\end{definition}

Se il primo documento rilevante si trova al rango 1, il punteggio vale $1.0$; se scende al rango 2 diviene $0.5$; se scende al rango 10 crolla a $0.1$. Qualora la graduatoria non contenga alcun documento utile entro la finestra considerata, si pone convenzionalmente $RR = 0$.

\FloatBarrier
\section{Oltre la Rilevanza Binaria: Rilevanza Graduata e NDCG}
\label{sec:rilevanza-graduata-ndcg}

Nelle applicazioni Web e commerciali, l'assunzione che un documento sia categoricamente rilevante ($1$) o non rilevante ($0$) si rivela eccessivamente semplicistica. 

Si consideri a titolo esemplificativo la query reale:
\begin{center}
\texttt{Query: ``information retrieval unipi''}
\end{center}
I risultati reperibili presentano sfumature qualitative nettamente distinte:
\begin{enumerate}
    \item \textbf{DidaWiki del Corso (Docente, CFU, programma didattico, dispense ufficiali):} Risponde integralmente e perfettamente all'intento dell'utente $\implies$ \textbf{Altamente Rilevante ($r = 3$)}.
    \item \textbf{Pagina Valutami UniPi (Scheda descrittiva e requisiti istituzionali d'esame):} Fonte autorevole e attinente, ma priva di materiale didattico operativo $\implies$ \textbf{Rilevante ($r = 2$)}.
    \item \textbf{Guida Generale del Corso di Laurea Magistrale:} Cita brevemente il corso nel piano di studi $\implies$ \textbf{Marginalmente Rilevante ($r = 1$)}.
    \item \textbf{Articoli di cronaca generale o pagine esterne slegate:} $\implies$ \textbf{Irrilevante ($r = 0$)}.
\end{enumerate}

Tutti i primi tre documenti presentano un legame tematico con la query, ma l'utente percepisce una soddisfazione radicalmente diversa. Per modellare questa granularit\`a, la letteratura adotta scale discrete di giudizio a pi\`u livelli (tipicamente $\mathcal{S} = \{0, 1, 2, 3\}$ o $\{0, 1, 2, 3, 4\}$).

\FloatBarrier
\subsection{I Fondamenti del Discounted Cumulative Gain (DCG)}
\label{subsec:fondamenti-dcg}

La formulazione del \textbf{Discounted Cumulative Gain} (DCG), introdotta da J\"arvelin e Kek\"al\"ainen~\cite{jarvelin2002}, poggia su due assunzioni comportamentali cardine:
\begin{enumerate}
    \item I documenti altamente rilevanti arrecano all'utente un'utilit\`a informativa (\textit{gain}) enormemente superiore rispetto a documenti solo marginalmente pertinenti.
    \item Minore \`e la posizione ordinale (rango) di un documento nella SERP, minore \`e la probabilit\`a che l'utente lo esamini; pertanto, il guadagno apportato deve essere progressivamente scontato (\textit{discounted}) verso il basso con il crescere del rango.
\end{enumerate}

\begin{definition}[Cumulative Gain a Rango $p$ ($CG_p$)]
Dato il vettore delle rilevanze graduate $\mathbf{r} = [r_1, r_2, \dots, r_p]$ fino al rango di taglio $p$, il Cumulative Gain \`e la somma lineare dei singoli guadagni:
\begin{equation}
CG_p = \sum_{i=1}^p r_i
\end{equation}
\textit{Limite strutturale:} $CG_p$ \`e del tutto insensibile all'ordinamento posizionale: scambiare un documento con $r=3$ tra il primo e il decimo posto non altera minimamente il valore di $CG_p$.
\end{definition}

\FloatBarrier
\subsection{Formulazione Standard vs. Formulazione Esponenziale}
\label{subsec:dcg-standard-vs-exp}

Per modellare il decadimento dell'attenzione visiva dell'utente, il DCG introduce un fattore di sconto posizionale:

\begin{definition}[DCG Standard (J\"arvelin e Kek\"al\"ainen)]
Applica uno sconto logaritmico in base al rango $i$:
\begin{equation}
DCG_p = r_1 + \sum_{i=2}^p \frac{r_i}{\log_2(i)}
\label{eq:dcg-standard}
\end{equation}
\end{definition}

\begin{property}[Significato Fisico dello Sconto Logaritmico]
Il divisore $\log_2(i)$ modella la progressione dello sguardo lungo la pagina:
\begin{itemize}
    \item Al rango $i = 4$, il fattore vale $\frac{1}{\log_2(4)} = \frac{1}{2}$: il valore del documento si dimezza rispetto alla prima posizione.
    \item Al rango $i = 8$, il fattore scende a $\frac{1}{\log_2(8)} = \frac{1}{3}$.
    \item Al rango $i = 16$, il fattore si riduce a $\frac{1}{\log_2(16)} = \frac{1}{4}$.
\end{itemize}
\end{property}

Nei motori di ricerca commerciali Web (Google, Microsoft Bing) si adotta prevalentemente la \textbf{formulazione esponenziale}:
\begin{equation}
DCG_p^{\text{exp}} = \sum_{i=1}^p \frac{2^{r_i} - 1}{\log_2(i + 1)}
\label{eq:dcg-exp}
\end{equation}

\begin{noteblock}{Differenze Chiave tra le Formulazioni di DCG}
\begin{enumerate}
    \item \textbf{Gestione del primo rango:} Nell'Equazione~\eqref{eq:dcg-standard} il termine $r_1$ \`e scorporato per evitare $\log_2(1) = 0$ a denominatore, facendo partire la sommatoria da $i=2$. Nell'Equazione~\eqref{eq:dcg-exp} l'indice parte direttamente da $i=1$ adottando $\log_2(i+1)$ ($\log_2(1+1) = \log_2(2) = 1$).
    \item \textbf{Guadagno non lineare:} Il numeratore $2^{r_i} - 1$ dilata esponenzialmente il peso dei documenti eccellenti: un documento con $r=3$ conferisce un guadagno $2^3 - 1 = 7$, mentre un documento con $r=1$ vale solo $2^1 - 1 = 1$. Questo premia in modo aggressivo il recupero dei documenti a massima pertinenza ai primissimi ranghi.
\end{enumerate}
\end{noteblock}

\FloatBarrier
\subsection{Normalized Discounted Cumulative Gain (NDCG)}
\label{subsec:normalized-dcg}

Poiché query diverse presentano un numero variabile di documenti rilevanti (e con differenti gradi di pertinenza massima), il valore assoluto di $DCG_p$ non \`e confrontabile né mediabile tra query eterogenee. 

Si definisce l'\textbf{Ideal DCG} ($IDCG_p$) come il valore di DCG calcolato sulla graduatoria ideale (\textit{perfect ranking}), ottenuta ordinando tutti i documenti della collezione in senso decrescente di rilevanza:

\begin{definition}[Normalized Discounted Cumulative Gain]
\begin{equation}
NDCG_p = \frac{DCG_p}{IDCG_p} \quad \in [0, 1]
\end{equation}
dove:
\begin{equation}
IDCG_p = r_1^* + \sum_{i=2}^p \frac{r_i^*}{\log_2(i)} \quad \text{con } r_1^* \ge r_2^* \ge \dots \ge r_p^*
\end{equation}
\end{definition}

La normalizzazione garantisce che un ordinamento perfetto consegua $NDCG = 1.0$ per qualsiasi query, rendendo la metrica idonea alla media macroscopica su ampie suite di benchmark.

\FloatBarrier
\subsection{Esempio Numerico Completo: Traccia a 10 Documenti}
\label{subsec:esempio-numerico-dcg-ndcg}

Si consideri la seguente graduatoria di 10 documenti, valutata su scala $r \in \{0, 1, 2, 3\}$:
\begin{equation}
\mathbf{r} = [3, \, 2, \, 3, \, 0, \, 0, \, 1, \, 2, \, 2, \, 3, \, 0]
\end{equation}
Ordinando i valori di rilevanza in senso non crescente si ottiene la graduatoria ideale:
\begin{equation}
\mathbf{r}^* = [3, \, 3, \, 3, \, 2, \, 2, \, 2, \, 1, \, 0, \, 0, \, 0]
\end{equation}

La Tabella~\ref{tab:ndcg-step-by-step} riporta l'evoluzione analitica dettagliata di ogni singolo termine posizionale secondo la formulazione standard.

\begin{table}[H]
\centering
\small
\resizebox{\columnwidth}{!}{%
\begin{tabular}{c c l c >{\bfseries}c c c >{\bfseries}c >{\bfseries}c}
\toprule
$i$ & $r_i$ & Divisore Sconto & Guadagno Scontato & $DCG@i$ & $r_i^*$ & Guadagno Ideale & $IDCG@i$ & $NDCG@i$ \\
\midrule
1  & 3 & $1.000$              & 3.00 & 3.00 & 3 & 3.00 & 3.00  & \textbf{1.00} \\
2  & 2 & $\log_2(2) = 1.000$  & 2.00 & 5.00 & 3 & 3.00 & 6.00  & \textbf{0.83} \\
3  & 3 & $\log_2(3) \approx 1.585$ & 1.89 & 6.89 & 3 & 1.89 & 7.89  & \textbf{0.87} \\
4  & 0 & $\log_2(4) = 2.000$  & 0.00 & 6.89 & 2 & 1.00 & 8.89  & \textbf{0.76} \\
5  & 0 & $\log_2(5) \approx 2.322$ & 0.00 & 6.89 & 2 & 0.86 & 9.75  & \textbf{0.71} \\
6  & 1 & $\log_2(6) \approx 2.585$ & 0.39 & 7.28 & 2 & 0.77 & 10.52 & \textbf{0.69} \\
7  & 2 & $\log_2(7) \approx 2.807$ & 0.71 & 7.99 & 1 & 0.36 & 10.88 & \textbf{0.73} \\
8  & 2 & $\log_2(8) = 3.000$  & 0.67 & 8.66 & 0 & 0.00 & 10.88 & \textbf{0.80} \\
9  & 3 & $\log_2(9) \approx 3.170$ & 0.95 & 9.61 & 0 & 0.00 & 10.88 & \textbf{0.88} \\
10 & 0 & $\log_2(10) \approx 3.322$ & 0.00 & 9.61 & 0 & 0.00 & 10.88 & \textbf{0.88} \\
\bottomrule
\end{tabular}%
}
\caption{Computazione analitica passo-passo di $DCG@i$, $IDCG@i$ e $NDCG@i$ per i 10 documenti dell'esempio canonico.}
\label{tab:ndcg-step-by-step}
\end{table}

\FloatBarrier
\subsection{Confronto tra Funzioni di Ranking Alternative}
\label{subsec:confronto-funzioni-ranking}

Per comprendere la sensibilit\`a di NDCG alle inversioni di rango, si analizzi un caso di studio con 4 documenti $\{d_1, d_2, d_3, d_4\}$ aventi gradi di rilevanza:
\begin{equation}
r(d_1) = 0, \quad r(d_2) = 1, \quad r(d_3) = 2, \quad r(d_4) = 2
\end{equation}

Si pongano a confronto due funzioni di ranking concorrenti con la graduatoria di riferimento (\textit{Ground Truth}):

\begin{itemize}
    \item \textbf{Ground Truth ($GT$):} Ordinamento $\langle d_4, d_3, d_2, d_1 \rangle$ con vettore $\mathbf{r}^* = [2, 2, 1, 0]$:
    \begin{align}
    IDCG_4 &= 2 + \frac{2}{\log_2(2)} + \frac{1}{\log_2(3)} + 0 \nonumber \\
    &= 2 + 2 + \frac{1}{1.5850} \approx 4 + 0.6309 = 4.6309 \nonumber \\
    NDCG_{GT} &= \frac{4.6309}{4.6309} = \mathbf{1.00}
    \end{align}
    \item \textbf{Ranking Function 1 ($RF_1$):} Ordinamento $\langle d_3, d_4, d_2, d_1 \rangle$ con vettore $\mathbf{r} = [2, 2, 1, 0]$.
    Sebbene $d_3$ e $d_4$ siano scambiati tra loro, poiché possiedono esattamente lo stesso grado di rilevanza ($r=2$), l'inversione non genera alcuna penalizzazione:
    \begin{align}
    DCG_{RF_1} &= 2 + \frac{2}{\log_2(2)} + \frac{1}{\log_2(3)} = 4.6309 \nonumber \\
    NDCG_{RF_1} &= \frac{4.6309}{4.6309} = \mathbf{1.00}
    \end{align}
    \item \textbf{Ranking Function 2 ($RF_2$):} Ordinamento $\langle d_3, d_2, d_4, d_1 \rangle$ con vettore $\mathbf{r} = [2, 1, 2, 0]$.
    In questo ordinamento, il documento subottimale $d_2$ ($r=1$) scavalca indebitamente $d_4$ ($r=2$) alla seconda posizione:
    \begin{align}
    DCG_{RF_2} &= 2 + \frac{1}{\log_2(2)} + \frac{2}{\log_2(3)} + 0 \nonumber \\
    &= 2 + 1 + \frac{2}{1.5850} \approx 3 + 1.2618 = 4.2618 \nonumber \\
    NDCG_{RF_2} &= \frac{DCG_{RF_2}}{IDCG_4} = \frac{4.2618}{4.6309} \approx \mathbf{0.9203}
    \end{align}
\end{itemize}

Questo esempio dimostra con rigore la propriet\`a fondamentale di NDCG: la funzione punisce proporzionalmente le discrepanze qualitative, distinguendo gli scambi neutri tra elementi paritari dalle inversioni penalizzanti tra documenti di rilevanza eterogenea.

\FloatBarrier
\subsection{Quadro Sinottico Comparativo delle Metriche di Valutazione}
\label{subsec:quadro-sinottico-metriche}

La Tabella~\ref{tab:sinottico-metriche} riassume sinotticamente l'intero panorama delle metriche esaminate nel capitolo, evidenziandone la natura matematica, il focus operativo e gli ambiti d'uso raccomandati.

\begin{table}[H]
\centering
\small
\resizebox{\columnwidth}{!}{%
\begin{tabular}{l l l l l}
\toprule
\textbf{Metrica} & \textbf{Rilevanza} & \textbf{Focus Principale} & \textbf{Punto di Forza} & \textbf{Limite Strutturale} \\
\midrule
\textbf{Precision ($P$)} & Binaria & Purezza insieme & Intuitiva, definita su insiemi & Ignora richiamo e ranking \\
\textbf{Recall ($R$)} & Binaria & Completezza estrazione & Misura copertura assoluta & Ignora rumore ($FP$) e ranking \\
\textbf{$F_1$-Score} & Binaria & Equilibrio armonico & Punisce forti asimmetrie & Ignora la graduatoria \\
\textbf{$F_\beta$-Score} & Binaria & Ponderazione flessibile & Regolabile per dominio & Ignora la graduatoria \\
\midrule
\textbf{P@k} & Binaria & Qualit\`a top-$k$ & Rispecchia la prima SERP & Non media bene se $R$ varia \\
\textbf{Recall@k} & Binaria & Frazione $R$ nei top-$k$ & Valuta richiamo locale & Dipende fortemente da $k$ \\
\textbf{AP} & Binaria & Area curva PR singola query & Sintesi globale dell'ordinamento & Complessa da calcolare su Web \\
\textbf{MAP} & Binaria & Macro-media su query suite & Tutela query di coda lunga & Assume rilevanza binaria \\
\textbf{MAP@K} & Binaria & Top-$K$ su Web e RecSys & Elimina distorsione di $P@k$ & Tronca la coda dei risultati \\
\textbf{MRR} & Binaria & Primo risultato utile & Ideale per factoid/known-item & Ignora tutti i risultati successivi \\
\midrule
\textbf{CG} & Graduata & Somma guadagni cumulativi & Supporta scale multi-livello & Cieca rispetto all'ordinamento \\
\textbf{DCG} & Graduata & Sconto posizionale logaritmico & Modella attenzione visiva reale & Non normalizzato tra query \\
\textbf{NDCG} & Graduata & Normalizzazione $DCG/IDCG$ & Standard assoluto ranking Web & Sensibile alla scelta del log/gain \\
\bottomrule
\end{tabular}%
}
\caption{Quadro sinottico comparativo delle metriche di valutazione per l'Information Retrieval.}
\label{tab:sinottico-metriche}
\end{table}
